\begin{ledgroupsized}[r]{120mm}
\footnotesize
\pstart
\noindent\textbf{Überlieferung:}
\pend
\end{ledgroupsized}
\begin{ledgroupsized}[r]{114mm}
\footnotesize
\pstart \parindent -6mm
\makebox[6mm][l]{\textit{L}}%
Konzept:
LH~XXXVII, 5~Bl.~183\textendash184. 
Ein Bogen~4\textsuperscript{o};
\href{\Scan{LH_37_05_184wz}}%
{Gegenmarke}
eines Wasserzeichens im Falz; %
Papiererhaltungsmaßnahmen;
alle Ränder beschnitten.
Vier Seiten.
Auf Bl.~184~r\textsuperscript{o}, senkrecht zur Schreibrichtung und von Leibniz überschrieben, ein Satz von der Hand 
\indexN{\textls{Hodann}, Johann Friedrich 1674\textendash1745}J.\,F.\,Hodanns:
\textit{Et vero de profectis in Graeciam %
loquuntur cippi, hominibus haud dubie Christia}\lbrack/\rbrack\textit{nis erecti,}
\textit{quos ipse
Cl.~\indexN{\textls{Peringer} (Perinskiöld, Perinskioldius), Johan 1654\textendash1720}%
Perinskioldius exhibet}.
Die Grabsteine in Griechenland gefallener gotischer Krieger sind Gegenstand von
\indexN{\textls{Peringer} (Perinskiöld, Perinskioldius), Johan 1654\textendash1720}%
J.\,Perinskiölds 
umfangreichen \glqq Annotationes\grqq\ zu seiner Ausgabe von
\indexN{\textls{Cochlaeus}, Johann 1479\textendash1552}%
\textsc{J.\,Cochlaeus}, 
\cite{02052}\textit{Vita Theoderici}, Stockholm 1699 (dort S.\;233\textendash614).
In diesem Zusammenhang siehe Leibnizens Kommentare über
\indexN{\textls{Peringer} (Perinskiöld, Perinskioldius), Johan 1654\textendash1720}%
Perinskiölds Thesen in seinem
\cite{02051}Brief an 
\indexN{\textls{Sparwenfeld}, Johan Gabriel 1655\textendash1727}%
J.\,G.~Sparwenfeld
vom 10.\ März 1702 (\textit{LSB} I, 20, S.\;838.1\textendash9.)
\pend
\end{ledgroupsized}
\vspace{5mm}
\begin{ledgroup}
\footnotesize
\pstart
\noindent%
\textbf{Datierungsgründe:}
Der von der Hand 
\indexN{\textls{Hodann}, Johann Friedrich 1674\textendash1745}Johann Friedrich Hodanns 
stammende, von Leibniz überschriebene Satz auf Bl.~184~r\textsuperscript{o} 
gibt Auskunft über die Entstehung des vorliegenden Konzeptes.
Hodann war
von Juni 1703 bis März 1705 und  von Februar 1706 bis Ende 1714
für Leibniz als Schreiber tätig.
Daraus ergibt sich für N.\;\ref{60354}
unter der Annahme einer zeitnahen Wiederverwendung des sonst unbeschriebenen Blattes durch Leibniz
die angegebene Datierungsspanne.
\pend 
\end{ledgroup}
\selectlanguage{latin}
\vspace{8mm}
\pstart%
\normalsize%
\noindent%
\edtext{\href{\Scan{LH_37_05_183r}}%
{\lbrack183~r\textsuperscript{o}\rbrack}\ Cum}{%
\lemma{\href{\Scan{LH_37_05_183r}}%
{\lbrack183~r\textsuperscript{o}\rbrack}}\Bfootnote{%
\textit{(1)}~In reddendis 
\textit{(a)}~legibus rationum 
\textit{(b)}~rationibus legum motus, duae Hypotheses 
\textit{(2)}~Cum~\textit{L}}} 
duo corpora directe et centraliter concurrunt,\indexS{concursus directus et centralis duorum} 
derivari possunt phaenomena\indexS{phaenomena} ex his duobus, primum si 
ponamus duo corpora concurrentia\indexS{corpora celeritatibus reciprocis concurrentia} celeritatibus ad 
magnitudines reciprocis\indexS{celeritas ad magnitudinem reciproca}
redire ea vi\indexS{vis} qua venere, secundo 
si ponamus idem evenire quacunque compositione motuum\indexS{compositio motuum}
\edtext{assumta, adeoque}{\lemma{assumta,}\Bfootnote{%
\textit{(1)}~ita n 
\textit{(2)}~seu 
\textit{(3)}~adeoque~\textit{L}}} 
communem motum\indexS{motus communis} duobus 
\edtext{corporibus supponendo}{\lemma{corporibus}\Bfootnote{%
\textit{(1)}~additum vel det 
\textit{(2)}~supponendo~\textit{L}}} 
\edtext{adjunctum, nihil mutare in motu proprio\indexS{motus proprius} ex concursu.}{\lemma{adjunctum,}\Bfootnote{%
\textit{(1)}~nihil mutare in concursibus, mod 
\textit{(2)}~nihil mutare in motu 
\textit{(a)}~speciali 
\textit{(b)}~proprio ex concursu.~\textit{L}}}
\pend\pstart
Prior Hypothesis in corporibus gravibus\indexS{corpus grave} demonstrari potest 
reducendo ad aliquod quasi absurdum, nempe ad motum perpetuum 
\edtext{mechanicum.\indexS{motus perpetuus mechanicus}}{\lemma{}\Bfootnote{mechanicum \textit{erg.~L}}} 
Porro 
\edtext{hujus}{\lemma{}\Bfootnote{hujus \textit{erg.~L}}} 
motus perpetui\indexS{motus perpetuus} absurditas logica necessitate non est evicta. Sed tantum duo certa 
\edtext{sunt\lbrack:\rbrack\ primum}{\lemma{sunt}\Bfootnote{%
\textit{(1)}~omnia 
\textit{(a)}~mechani 
\textit{(b)}~ten
\textit{(2)}~primum~\textit{L}}} 
ab 
\edtext{experimento, omnia tentamenta}{\lemma{experimento,}\Bfootnote{%
\textit{(1)}~omnem 
\textit{(2)}~omnia tentamenta~\textit{L}}} 
motus perpetui\indexS{tentamentum motus perpetui} vana reperiri; secundo a 
\edtext{ratione metaphysica,}{\lemma{ratione}\Bfootnote{%
\textit{(1)}~, quod 
\textit{(2)}~metaphysica~\textit{L}}} 
sed probabili magis quam rigorose demonstrata, 
quod effectus plenus\indexS{effectus plenus}\indexS{effectus plenus non fortior causa} non debeat esse fortior causa.\indexS{causa} 
Equidem rigorose demonstrari potest in isorrhopicis\indexS{isorrhopica}, seu 
\edtext{staticis\indexS{statica} ubi}{\lemma{staticis}\Bfootnote{%
\textit{(1)}~non dari posse 
\textit{(2)}~ubi~\textit{L}}} 
de vi mortua\indexS{vis mortua} agitur, non dari posse motum perpetuum.\indexS{motus perpetuus}
Sed hic de vi 
\edtext{viva\indexS{vis viva} nulla}{\lemma{viva}\Bfootnote{%
\textit{(1)}~non 
\textit{(2)}~nulla~\textit{L}}} 
hujusmodi suppetit demonstratio, 
\edtext{in gravibus,\indexS{grave}}{\lemma{}\Bfootnote{in gravibus, \textit{erg.~L}}} 
eventus tamen non minus 
\edlabel{37_05_183-184_1a}%
\edtext{}{%
{\xxref%
{37_05_183-184_1a}{37_05_183-184_1b}}%
\lemma{respondet.}\Bfootnote{%
\textit{(1)}~Sed 
\textit{(2)}~Equidem~\textit{L}}}%
respondet.
\pend
\pstart
Equidem%
\edlabel{37_05_183-184_1b}
cum duo corpora horizonti\indexS{horizon} sunt parallela,\indexS{corpora horizonti parallela} gravitatis\indexS{gravitas} nulla, aut pene nulla actio\indexS{actio} est, 
et tamen res docet observari hanc reciprocationem materiae et 
\edtext{velocitatis\indexS{reciprocatio materiae et velocitatis} etiam}{\lemma{velocitatis}\Bfootnote{%
\textit{(1)}~etiam in ip 
\textit{(2)}~non ta 
\textit{(3)}~etiam~\textit{L}}} 
in infinite parva vi,\indexS{vis infinite parva} seu solicitatione,\indexS{solicitatio}
si ab Elastro\indexS{Elastrum} omnia deducamus unde oritur idem deinde in vi 
\edlabel{37_05_183-184_2a}%
\edtext{}{%
{\xxref%
{37_05_183-184_2a}{37_05_183-184_2b}}%
\lemma{viva.}\Bfootnote{%
\textit{(1)}~Sed 
\textit{(2)}~Idque in
\textit{(3)}~Sup 
\textit{(4)}~Satis~\textit{L}}}%
viva.\indexS{vis viva} %
\pend
\pstart
Satis%
\edlabel{37_05_183-184_2b}
ergo esset si 
in vi mortua\indexS{vis mortua} seu solicitatione\indexS{solicitatio}
demonstrari 
\edtext{posset haec}{\lemma{posset}\Bfootnote{%
\textit{(1)}~solicitatio 
\textit{(2)}~haec~\textit{L}}}
reciprocatio.\indexS{reciprocatio materiae et velocitatis}
\edtext{Sane in gravitate\indexS{gravitas} eam demonstrat Archimedes,\indexN{\textls{Archimedes} 287\textendash212 v. Chr.} ex eo 
\edtext{in effectu\lbrack:\rbrack}{\lemma{}\Bfootnote{in effectu \textit{erg.~L}}} 
quod majus 
\edtext{plures habet}{\lemma{plures}\Bfootnote{%
\textit{(1)}~recipit i 
\textit{(2)}~habet~\textit{L}}} 
in se descendendi causas, adeoque praevalet, pro magnitudinis ratione etc.}{%
\lemma{Sane \lbrack...\rbrack\ etc.}%
\Cfootnote{%
\indexN{\textls{Archimedes} 287\textendash212 v. Chr.}\textsc{Archimedes}, 
\cite{01010}\textit{De planorum aequilibris libri duo}, lib.~I, prop.~II.}}
\edtext{Sed cum}{\lemma{Sed}\Bfootnote{%
\textit{(1)}~in 
\textit{(2)}~cum~\textit{L}}} 
\edtext{gravitate\indexS{gravitas} 
\edtext{s\lbrack e\rbrack posita,}{%
\lemma{}%
\Bfootnote{%
supposita \textit{L ändert Hrsg.}}}
nullum tale apparet principium et tamen gravitate\indexS{gravitas} 
\edtext{s\lbrack e\rbrack posita}{%
\lemma{}%
\Bfootnote{%
supposita \textit{L ändert Hrsg.}}}%
}{%
\lemma{gravitate \lbrack...\rbrack\  s\lbrack e\rbrack posita}%
\Cfootnote{%
Die in \textit{L} zweimal vorkommende Wendung \glqq gravitate supposita\grqq\ ist mit Leibnizens Anspruch, den Reziprozitätssatz (im Gegensatz zur archimedischen Demonstration) unabhängig von der Schwere der Körper zu beweisen, nicht vereinbar.
Der Wortlaut wird in Anlehnung an die spätere Stelle \refpassage{37_05_183-184_4a}{37_05_183-184_4b} (\glqq ut lex staticorum seu aequilibrii etiam seposita gravitate succedat\grqq) emendiert.}}
idem habere locum ostendit res ipsa, nam corpus
\edtext{etiam grave\indexS{corpus grave}}{\lemma{}\Bfootnote{etiam grave \textit{erg.~L}}} 
ascendens impetu\indexS{impetus} quodam, alterique descendenti occurrens 
deprimere conanti renititur, perinde ac si grave non esset. Atque id quidem fieri debere ostenditur per 
consequentiam, ut evitetur motus perpetuus mechanicus.\indexS{motus perpetuus mechanicus}
Sed hanc causam finalem\indexS{causa finalis} non apparet, qua
\href{\Scan{LH_37_05_183v}}%
{\lbrack183~v\textsuperscript{o}\rbrack}\ efficiente\indexS{causa efficiens} obtineat 
\edlabel{37_05_183-184_3a}%
\edtext{}{%
{\xxref%
{37_05_183-184_3a}{37_05_183-184_3b}}%
\lemma{natura.}\Bfootnote{%
\textit{(1)}~Equidem commentum \bars{quoddam \textit{streicht Hrsg.}} excogitari 
\textit{(2)}~Sane~\textit{L}}}%
natura.
\pend
\pstart
Sane%
\edlabel{37_05_183-184_3b}
si fingamus omne 
\edtext{corpus ipsa}{\lemma{corpus}\Bfootnote{%
\textit{(1)}~\bars{vi \textit{streicht Hrsg.}} quadam 
\textit{(2)}~ipsa vi~\textit{L}}} 
vi\indexS{vis} impellentium vicinorum\indexS{vis corporum vicinorum impellentium} conspirante in loco retineri, et ita quodammodo 
respectu loci\indexS{locus corporis} in quo est, dum in eo est\lbrack,\rbrack\ grave esse, res procedit; 
et ratio redditur hujus majoris resistentiae\indexS{resistentia ad motum} 
\edtext{ad motum}{\lemma{}\Bfootnote{ad motum \textit{erg.~L}}} 
in corpore majore\indexS{corpus majus}. Sed cum ipsa haec 
\edtext{corpora minora\indexS{corpus minus} quorum}{\lemma{corpora}\Bfootnote{%
\textit{(1)}~quae retinent 
\textit{(2)}~minora quorum~\textit{L}}} 
motu et velut flatu corpus majus in loco\indexS{locus corporis} retinetur credibile sit subjecta esse iisdem legibus motus;\indexS{lex motus}
jam patet causam mechanicam\indexS{causa mechanica} ex efficientibus\indexS{causa efficiens} reddi non posse. 
\pend
\pstart
Idem est si 
\edtext{fingamus materiam}{\lemma{fingamus}\Bfootnote{%
\textit{(1)}~corpora a se invi 
\textit{(2)}~materiam~\textit{L}}} 
distributam in corpora aequalia sese non tangentia, quibus 
interfluat fluidum tenue,\indexS{fluidum tenue} quod ipsa ad se invicem appropinquare non sinat, 
sed a se invicem repellat, hac fictione\indexS{fictio} 
\edtext{explicari potest hoc \greek{f}aenomenon quod}{\lemma{explicari}\Bfootnote{%
\textit{(1)}~possunt haec \greek{f}aenomena de causa cur 
\textit{(2)}~potest hoc \greek{f}aenomenon quod~\textit{L}}} 
majus corpus\indexS{corpus majus} tanto tardius moveatur, sed 
\edtext{quia fluidum\indexS{fluidum} ipsum non potest esse
\edtext{purum, per}{%
\lemma{purum,}%
\Bfootnote{%
\textit{(1)}~ex %
\textit{(2)}~per~\textit{L}%
}}
demonstrationes nostras metaphysicas,\indexS{demonstratio metaphysica}}{%
\lemma{fluidum \bracks{...} metaphysicas}%
\Cfootnote{%
Siehe bspw. \cite{02557}\textit{Specimen Dynamicum} \bracks{Zweiter Teil}, \cite{01043}\textit{LMG}~VI, S.\;246-254, bes. S.\;248\,f.}}
sed apparens tantum, 
rursusque in particulas\indexS{particulae} resolvitur, atque ita in 
\edtext{infinitum, quae}{%
\lemma{infinitum,}%
\Bfootnote{%
\textit{(1)}~patet \bars{semper \textit{streicht Hrsg.}}
\textit{(a)}~su %
\textit{(b)}~suppo %
\textit{(2)}~quae~\textit{L}}}
et ipsae easdem subeant leges, patet semper id 
\edtext{quod explicuimus}{%
\lemma{quod}%
\Bfootnote{%
\textit{(1)}~in %
\textit{(2)}~explicuimus~\textit{L}}}
rursus explicari debere, atque adeo altius aliquid 
\edtext{subesse: nec}{\lemma{subesse:}\Bfootnote{%
\textit{(1)}~ipsam 
\textit{(2)}~nec~\textit{L}}} 
video quid aliud esse possit, 
\edtext{quam ut dicamus esse}{\lemma{quam}\Bfootnote{%
\textit{(1)}~esse 
\textit{(2)}~ut dicamus esse~\textit{L}}} 
naturalem corporum inertiam,\indexS{inertia naturalis corporum} seu hoc ipsum 
consilium naturae\indexS{consilium naturae} omnia ita explicandi, ut fictio illa succedat 
\edtext{de partibus}{\lemma{de}\Bfootnote{%
\textit{(1)}~interfluente et 
\textit{(2)}~partibus~\textit{L}}} 
aequalibus et fluido interfluente.\indexS{fluidum interfluens}
\pend
\pstart
Succedit etiam fictio, si ponamus corpus ad locum\indexS{locus corporis} suum gravitare, 
dum in eo est\lbrack,\rbrack\ tanquam 
\edtext{impulsu corporis}{\lemma{impulsu}\Bfootnote{%
\textit{(1)}~extraneorum, etsi 
\textit{(2)}~rursus de ipsis c 
\textit{(3)}~corporis~\textit{L}}} 
fluidi ambientis;\indexS{fluidum ambiens} sed rursus hoc fluidum\indexS{fluidum} 
pares leges\indexS{lex} subire debet. Itaque 
\edtext{naturae consilium\indexS{consilium naturae}}{\lemma{naturae}\Bfootnote{%
\textit{(1)}~aliquid 
\textit{(2)}~consilium~\textit{L}}} 
fuit efficiendi ut succedant hae fictiones;\indexS{fictio} idque ipsum dici potest 
inertiam corporum naturalem\indexS{inertia naturalis},
\edlabel{37_05_183-184_4a}%
ut lex staticorum\indexS{lex staticorum} seu aequilibrii\indexS{lex aequilibrii}
etiam seposita gravitate\indexS{gravitas} succedat.%
\edlabel{37_05_183-184_4b}
Et est haec de successu fictionum\indexS{fictio}\indexS{successus fictionum} 
observatio pene eadem cum 
\href{\Scan{LH_37_05_184r}}%
{\lbrack184~r\textsuperscript{o}\rbrack}\
observatione 
\edtext{%
\edtext{regulae meae,}{%
\lemma{regulae}%
\Bfootnote{%
\textit{(1)}~de transitu %
\textit{(2)}~meae,~\textit{L}}}
quam appello legem continuitatis,\indexS{lex continuitatis}}{%
\lemma{regulae \lbrack...\rbrack\ continuitatis}%
\Cfootnote{%
Siehe \cite{02054}\textit{Extrait d'une Lettre de M.~L.\ sur un Principe Général} in 
\cite{02002}\textit{NRL}, Juli 1687, S.\;744\textendash753 
(\cite{02039}\textit{LPG} III, S.\;51\textendash55)
sowie 
\cite{02055}\textit{Principium quoddam generale}, 
\textit{LSB} VI, 4, N.\;371.}}	%
quae in applicatione ad \greek{f}ysica\indexS{physica} non demonstratur rigorose 
ut in Geometricis,\indexS{Geometrica} et tamen succedit, 
quod alias natura ordinis\indexS{ordo naturae} expers deprehendatur. 
Unde eadem prodire 
\edtext{debent sive corpus}{\lemma{debent}\Bfootnote{%
\textit{(1)}~ex lege 
\textit{(2)}~sive 
\textit{(a)}~corpora 
\textit{(b)}~corpus~\textit{L}}} 
in corpora plura 
\edtext{incurrat}{\lemma{}\Bfootnote{incurrat \textit{erg.~L}}} 
simul,\indexS{corpus in plura simul incurrens} sive 
\edtext{successive temporis}{\lemma{successive}\Bfootnote{%
\textit{(1)}~exiguo 
\textit{(2)}~temporis~\textit{L}}} 
intervallo infinite parvo.\indexS{intervallum infinite parvum}
\edtext{Quod principium an in \greek{f}ysicis succederet
Jungius\indexN{\textls{Jungius}, Joachim 1587\textendash1657}
in schedis Ms. logicis aliquando dubitavit.%
}{\lemma{Quod \lbrack...\rbrack\ dubitavit}%
\Cfootnote{%
Die Anspielung konnte noch nicht aufgelöst werden. %
}} 
\pend
\pstart
Quod etiam natura non patitur se decipi, seu motum perpetuum\indexS{motus perpetuus} 
arte sibi obtrudi, seu quod non 
\edtext{vult effectum}{\lemma{vult}\Bfootnote{%
\textit{(1)}~causam 
\textit{(2)}~effectum~\textit{L}}} 
causae praepollere,\indexS{effectus causae non praepollens} 
reducendum est ad aliquod intrinsecum inexistens. Et quidem 
\edlabel{37_05_183-184_5a}%
\edtext{}{%
\xxref{37_05_183-184_5a}{37_05_183-184_5b}%
\lemma{reperi \bracks{...} intelligi}%
\Cfootnote{Siehe bspw. \cite{02063}\textit{De ipsa natura}, \cite{01023}\textit{AE} 1698, S.\;427--440, bes. Art. (4)--(6), S.\;429--431 (\cite{02039}\textit{LPG} IV, S.\;504--516, bes. S.\;505--507).}}%
\edtext{hoc eo}{\lemma{hoc}\Bfootnote{%
\textit{(1)}~etiam 
\textit{(2)}~eo~\textit{L}}} 
reduci\lbrack,\rbrack\ ut eadem semper maneat quantitas 
\edtext{actionis,\indexS{quantitas actionis} seu}{\lemma{actionis,}\Bfootnote{%
\textit{(1)}~quod potest intelligi insitum 
\textit{(2)}~seu~\textit{L}}} 
eadem vis\indexS{vis}\lbrack,\rbrack\ quod
\edtext{potest insitum intelligi.\edlabel{37_05_183-184_5b} Sed utendo principio aequipollentis actionis vivae,\indexS{actio viva}\indexS{principium aequipollentis actionis vivae} non proceditur}{\lemma{potest}\Bfootnote{%
\textit{(1)}~esse 
\textit{(2)}~insitum intelligi.
\textit{(a)}~Et ex hoc principio optime demonstratur lex inertiae, quod 
\textit{(b)}~Sed ita res non \bars{reducitur \textit{streicht Hrsg.}} ad le 
\textit{(c)}~Sed princi 
\textit{(d)}~Sed ita 
\textit{(e)}~Nam 
\textit{(f)}~Sed utendo principio 
\textit{(aa)}~permanentis 
\textit{(bb)}~aequipollentis actionis vivae, 
\textit{(aaa)}~non observatur vis mo 
\textit{(bbb)}~proceditur
\textit{(ccc)}~non proceditur~\textit{L}}} 
per vim mortuam\indexS{vis mortua} seu viam aequilibrii\indexS{aequilibrium}
Verum enim vero dici 
\edtext{potest, ipsam}{\lemma{potest,}\Bfootnote{%
\textit{(1)}~in i 
\textit{(2)}~ipsam~\textit{L}}} 
vim mortuam\indexS{vis mortua} et vivam\indexS{vis viva} 
in eodem quasi temporis momento posse intelligi, 
ut in eodem puncto 
tam osculum\indexS{osculum} quam contactus,\indexS{contactus} seu tam directio\indexS{directio} quam flexus\indexS{flexus}
etsi se invicem se habeant, ut 
vis viva\indexS{vis viva} et mortua.\indexS{vis mortua}
Ut taceam vim mortuam\indexS{vis mortua} ipsam non prodire 
nisi ut consequentiam vivae,\indexS{vis viva} velut 
\edtext{ex circulatione\indexS{circulatio}}{\lemma{ex}\Bfootnote{%
\textit{(1)}~principio 
\textit{(2)}~circulatione~\textit{L}}}
oritur recessus a tangente.\indexS{recessus a tangente}%
Ergone dicemus legem de conservatione quantitatis actionum\indexS{quantitas actionis}\indexS{lex de conservatione quantitatis actionum} natura\indexS{natura} priorem esse.
\edtext{Sed contra, stat lex mea nova, de omni transitu per infinite parva.\indexS{lex de omni transitu per infinite parva}}{\lemma{}\Bfootnote{Sed \lbrack...\rbrack\ parva. \textit{erg.~L}}}
An in 
\edtext{talibus, cum}{\lemma{talibus,}\Bfootnote{%
\textit{(1)}~naturam 
\textit{(2)}~nih 
\textit{(3)}~nihi 
\textit{(4)}~cum~\textit{L}}} 
mera sint phaenomena\indexS{phaenomena}\lbrack,\rbrack\
\edtext{dicemus}{\lemma{}\Bfootnote{dicemus \textit{erg.~L}}} 
nihil esse prius vel posterius natura intendente, ut res utroque 
\edtext{modo concepta}{\lemma{modo}\Bfootnote{%
\textit{(1)}~succedat 
\textit{(2)}~concepta~\textit{L}}}
succedat. Generaliter dici potest vim\indexS{vis} manere, 
\edtext{seu effectus quantitatem\indexS{quantitas effectus} sive in mortua\indexS{vis mortua} sive viva vi\indexS{vis viva}, et}{\lemma{seu}\Bfootnote{%
\textit{(1)}~potestatem 
\textit{(2)}~effectus quantitatem
\textit{(a)}~,~\bars{et \textit{streicht Hrsg.}} quod aliud loco dimovet
\textit{(b)}~sive \lbrack...\rbrack\ et~\textit{L}}}
superare
\edtext{resistentiam\indexS{resistentia}\lbrack,\rbrack\ quod}{\lemma{resistentiam}\Bfootnote{%
\textit{(1)}~alterique vim dare 
\textit{(2)}~dum alteri 
\textit{(3)}~quod~\textit{L}}} 
alteri dat vim\indexS{vis}\lbrack,\rbrack\ et rationem esse cur succedant omnes fictiones.\indexS{successus fictionum}
\pend
\pstart Quod si jam veniamus ad alterum de compositione motuum,\indexS{compositio motuum} 
sequitur ex ea corpus quiescens\indexS{corpus quiescens} vim habere seu concipi posse 
ac si duos haberet motus 
\edtext{contrarios.\indexS{motus duo contrarii}
\href{\Scan{LH_37_05_184v}}%
{\lbrack184~v\textsuperscript{o}\rbrack}\
Loco compositionis motuum\indexS{compositio motuum} possumus 
uti principio}{%
\lemma{contrarios. \href{\Scan{LH_37_05_184v}}%
{\lbrack184~v\textsuperscript{o}\rbrack}}\Bfootnote{%
\textit{(1)}~Et 
\textit{(2)}~Ex 
\textit{(3)}~Dum principio illo utimur,
\textit{(4)}~Loco \lbrack...\rbrack\ quod~\textit{L}}} 
ictus 
\edtext{aequalis\indexS{principium ictus aequalis} dum}{\lemma{aequalis}\Bfootnote{%
\textit{(1)}~sit utrum corporum concurrat, non ideo 
\textit{(2)}~in 
\textit{(3)}~dum~\textit{L}}} 
eadem sit celeritas appropinquationis\indexS{celeritas appropinquationis}, 
quaecunque corpora concurrant, sed 
\edtext{hoc}{\lemma{}\Bfootnote{hoc \textit{erg.~L}}} 
ipsum supponit quiescentis resistentiam,\indexS{resistentia} 
aut dicendum est, neque motum\indexS{motus} neque quietem\indexS{quies} esse quicquam 
reale, sed tantum mutationem situs.\indexS{mutatio situs}
Verum ita nec ipsa quicquam reale erit, cum nulla possit intelligi duorum 
situs mutatio,\indexS{mutatio situs} quin simul determinatum esse oporteat, quodnam ex ipsis et quantum moveatur. Et ratio 
motus\indexS{ratio motus} sumi debet ex eo corpore quod cum alio concurrens 
\edtext{ideo}{\lemma{}\Bfootnote{ideo \textit{erg.~L}}} 
mutavit 
\edtext{directionem. Saltem}{\lemma{directionem.}\Bfootnote{%
\textit{(1)}~In eo enim revera facta \bars{est \textit{streicht Hrsg.}} %
\textit{(2)}~Saltem~\textit{L}}} 
enim in altero eorum necessario facta est mutatio.\indexS{mutatio}
Itaque in realibus principium est discernendi. Et tamen nullum est in resultantibus.
Cujus rei quae causa sit intelligi potest. Itaque rursus fictionis\indexS{fictio} succedentis
genus est, haec Hypothesis ut eadem 
\edtext{proveniant quomodocunque}{\lemma{proveniant}\Bfootnote{%
\textit{(1)}~sive 
\textit{(2)}~quomodocunque~\textit{L}}} 
producatur motus aliquis absolutus\indexS{motus absolutus} vel quies aliqua absoluta.\indexS{quies absoluta} 
Sive verb.\,gr. quies 
\edtext{sit realis}{\lemma{sit}\Bfootnote{%
\textit{(1)}~vera 
\textit{(2)}~realis~\textit{L}}} 
\edtext{sive motus}{\lemma{sive}\Bfootnote{%
\textit{(1)}~quies 
\textit{(2)}~motus~\textit{L}}} 
duo contrarii\indexS{motus duo contrarii} sint 
\edtext{compositi,\indexS{compositio duorum motuum contrariorum} ut corpore}{\lemma{compositi,}\Bfootnote{%
\textit{(1)}~ut 
\textit{(a)}~se 
\textit{(b)}~in navi 
\textit{(2)}~ut
\textit{(a)}~corpora
\textit{(b)}~corpore~\textit{L}}} 
contra navem\indexS{navis} in navi\indexS{navis} currente.\indexS{corpus contra navem in navi currens} 
Haec Hypothesis etiam mira est, et tamen succedit. %
\pend
\pstart
Quid si hanc adhibeamus fictionem, 
\edtext{ut corpora}{\lemma{ut}\Bfootnote{%
\textit{(1)}~corpus 
\textit{(2)}~corpora~\textit{L}}} 
duo concurrentia ponantur divisa in partes aequales infinite
\edtext{parvas,\indexS{partes infinite parvae}}{\lemma{}\Afootnote{\textit{Am Rand, neben} infinite:\enskip indefinite\indexS{partes indefinite parvae}}} 
et motus inaequalis etiam in motus aequales, sed quietibus 
\edtext{interruptos.\indexS{motus aequales quietibus interrupti} Quae}{\lemma{interruptos.}\Bfootnote{%
\textit{(1)}~Ita 
\textit{(2)}~Quae~\textit{L}}} 
fictio semper imminuta tandem desinat in statum 
\edtext{praesentem, motus continui\indexS{motus continuus} et}{\lemma{praesentem,}\Bfootnote{%
\textit{(1)}~dum tandem 
\textit{(2)}~motus \bars{continui \textit{erg.}} et~\textit{L}}} 
corporis continui.\indexS{corpus continuum} Oportet 
\edtext{autem supponere corpus}{\lemma{autem}\Bfootnote{%
\textit{(1)}~fingere co 
\textit{(2)}~supponere corpus~\textit{L}}} 
incurrens in corpus aequale quiescens\indexS{corpus in aequale quiescens incurrens}
suum ei motum dare, ejusque loco 
\edtext{quiescere. Videtur}{\lemma{quiescere}\Bfootnote{%
\textit{(1)}~, vel 
\textit{(2)}~. Videtur~\textit{L}}} 
duplici modo procedi posse in hac fictione, uno ut 
\edtext{corpora concurrant}{\lemma{corpora}\Bfootnote{%
\textit{(1)}~sibi occurrant 
\textit{(2)}~concurrant~\textit{L}}} 
tunc cum unum quiescit, idque alternis; ita non opus erit 
\edtext{motus inaequales}{\lemma{motus}\Bfootnote{%
\textit{(1)}~aequales 
\textit{(2)}~inaequales~\textit{L}}} 
reducere ad aequales; alter ut corpora concurrant inter movendum. Ubi rursus non opus erit 
motum inaequalem resolvi in 
\edtext{motus}{\lemma{}\Bfootnote{motus \textit{erg.~L}}} 
aequales interruptos,\indexS{motus aequales interrupti} si ponamus corpora inaequalium motuum 
in concursu\indexS{concursus} permutare motus.\indexS{permutatio motuum}
\pend
\pstart Concipiendum in momento concursus corpora tantum sibi contraria esse seu 
ictum\indexS{ictus} exercere vi motus proprii.\indexS{motus proprius}
\pend
